% Report for the NLP course project P6 The Apprentice Model
% Based on the Springer Nature article template

\documentclass[pdflatex,sn-mathphys-num]{sn-jnl}

\usepackage{graphicx}
\usepackage{amsmath,amssymb,amsfonts}
\usepackage{booktabs}
\usepackage{xcolor}
\usepackage{textcomp}
\usepackage{url}
\usepackage{listings}


\raggedbottom

\begin{document}

\title[The Apprentice Model for Toxicity Detection]{Distilling a BERT-base Teacher into a Compact Transformer for Toxic Comment Classification}

\author{\fnm{Vadim} \sur{Sokolov}}

\affil{\orgdiv{Master's Degree in Computer Science}, \orgname{Universit\`a degli Studi di Milano}, \orgaddress{\city{Milan}, \country{Italy}}}

\abstract{This project studies knowledge distillation for Natural Language Processing through an ``apprentice model''. 
The goal is to investigate if a compact student model can imitate the decision behaviour of a larger teacher model on a specific text classification task.
The selected task is toxic comment classification. 
The experimental pipeline compares a classical TF-IDF and Logistic Regression baseline, a supervised BERT-tiny student, a supervised BERT-base teacher trained on original dataset labels, and a BERT-tiny student distilled from BERT-base logits.
On the natural distribution test split, the TF-IDF baseline got an F1 score of 0.406 at the default threshold of 0.5.
The supervised BERT-base teacher obtains an F1 score of 0.620.
The supervised BERT-tiny baseline obtains F1 0.504, and the soft-label distilled BERT-tiny model with $T=2.0$ and $\alpha=0.7$ obtains F1 0.565.
This recovers 52.65\% of the F1 gap between supervised BERT-tiny and BERT-base while keeping the compact BERT-tiny model size and inference cost.}

\keywords{Natural Language Processing, Knowledge Distillation, Toxicity Detection, BERT, TF-IDF, Large Language Models, GPT}

\maketitle

\section*{Generative AI disclaimer}

Generative AI, specifically GPT-5.5, was used during the project implementation to validate the overall approach, suggest ideas for improving the experimental results, check the grammar of the report, create LaTeX tables, generate experiment configuration files, generate code documentation with comments, and help resolve Python package conflicts in the virtual environment.

\section{Introduction}\label{sec:introduction}

Large language models have demonstrated strong performance across many Natural Language Processing tasks, including classification, question answering, summarisation, and reasoning.
However, their size and computational cost make them difficult to utilize in constrained environments, especially when the target application requires limited resources.
Knowledge distillation resolves this problem by transferring knowledge from a larger teacher model to a smaller student model \cite{hinton2015distilling}.
Instead of trying to reproduce the full generality of the teacher, the student can be trained to imitate the teacher's behaviour on a specific task or domain.

This project explores this idea in the context of toxic comment classification. Toxicity detection is a relevant NLP task because it combines text classification, contextual language understanding, and ethical considerations related to bias, moderation, and disagreement.
The task is also appropriate for studying the tradeoff between model performance and efficiency. A large LLM may provide high quality judgments but it's slow and expensive to use.
Meanwhile, a compact model can be executed locally and quickly.

The project investigates if a compact Transformer model can be used as a lightweight student of a larger teacher model.
The selected student is BERT-tiny, a small Transformer encoder.
The final teacher is a finetuned BERT-base classifier.
In addition, a classical TF-IDF and Logistic Regression model is used as a baseline, connecting the project to traditional vector space methods covered in the course.

The relevant literature includes the original formulation of knowledge distillation by Hinton et al. \cite{hinton2015distilling}, where a student model learns from softened teacher output distributions and Transformer based language models such as BERT \cite{devlin2019bert}.
TinyBERT \cite{jiao2020tinybert} demonstrate that smaller Transformer models can preserve part of the performance of larger models, reducing computational cost. 
This project applies the same general motivation to a practical toxicity detection pipeline and evaluates the resulting trade offs in task performance, model size, and inference speed.

\section{Research question and methodology}\label{sec:methodology}

\subsection{Research question}

The main research question is:


To what extent can a compact Transformer model imitate a stronger BERT-base teacher for toxic comment classification.
What are the trade offs between predictive performance, model size, and inference speed?


This question is task specific. 
The goal is not to show that a small model can replace a general purpose LLM in reasoning or conversation. 
The project studies if the teacher's decision behaviour can be compressed into a smaller model for one domain specific classification problem.

\subsection{Problem definition}

Let $x_i$ be a text comment and let $y_i \in \{0,1\}$ be its original dataset label, where $0$ denotes not toxic and $1$ denotes toxic. 
The supervised baseline learns a classifier
\begin{equation}
    f_{\theta}(x_i) \rightarrow \hat{y}_i,
\end{equation}
using the original dataset labels. 


For the final distillation experiment, the teacher $T$ is a finetuned BERT-base classifier.
For each example, the teacher produces logits $z_T(x_i)$ and probabilities $p_T(y=0|x_i)$ and $p_T(y=1|x_i)$.
The resulting distillation dataset is
\begin{equation}
    \mathcal{D}_T = \{(x_i, y_i, z_T(x_i))\}_{i=1}^{N}.
\end{equation}

The student $S_{\phi}$ is a BERT-tiny classifier trained with soft-label temperature distillation.
The loss combines cross entropy against the original label with a KL-divergence term that matches the softened teacher and student distributions:

\begin{equation}
    \mathcal{L} = (1 - \alpha)\,\mathcal{L}_{CE}(y_i, p_S(x_i)) + \alpha \tau^2\,\mathcal{L}_{KL}(p_T^{(\tau)}(x_i), p_S^{(\tau)}(x_i)),
\end{equation}
where $\mathcal{L}_{CE}$ is cross entropy with hard labels, $\mathcal{L}_{KL}$ is the Kullback--Leibler divergence between teacher and student probability distributions, $p_T^{(\tau)}(x_i)=\operatorname{softmax}(z_T(x_i)/\tau)$, $p_S^{(\tau)}(x_i)=\operatorname{softmax}(z_S(x_i)/\tau)$, $\tau$ is the distillation temperature, and $\alpha$ controls the trade off between the two terms.
The final experiment uses logits from the finetuned BERT-base teacher with $\tau=2.0$ and $\alpha=0.7$.

\subsection{Pipeline overview}

The project pipeline consists of five stages:

\begin{enumerate}
    \item \textbf{Data preparation.} The dataset is loaded from Hugging Face and converted into reproducible train, validation, and test CSV splits. 
    The main supervised experiments use natural-distribution stratified splits that preserve the original class imbalance.
    \item \textbf{Baseline models.} A TF-IDF and Logistic Regression model is trained as a baseline. 
    A BERT-tiny model is then finetuned on the original labels as a supervised Transformer baseline.
    \item \textbf{Supervised local teacher.} A BERT-base model is finetuned on the natural distribution original labels.
    It is used to produce logits and probabilities for later soft-label distillation.
    \item \textbf{Distilled student.} Train a distilled BERT-tiny student from BERT-base teacher logits using soft-label temperature distillation. 
    Evaluate quality metrics and benchmark model size and inference speed.
\end{enumerate}

The implementation is organised as a Python package.
Separate modules are used for data loading, preprocessing, baseline models, Transformer training, teacher prompting, evaluation, and benchmarking. 

\section{Experimental results}\label{sec:results}

\subsection{Dataset and preprocessing}

The experiments use the \texttt{SetFit/toxic\_conversations} dataset from Hugging Face.
The dataset contains comments labelled as toxic or not toxic and the columns \texttt{text}, \texttt{label}, and \texttt{label\_text}. 
The full dataset contains more than 1.7 million training examples and 50,000 test examples.
The original distribution is highly imbalanced, with around 8\% toxic examples. 
The main original-label experiments use natural distribution stratified splits.
The training and validation splits are sampled from the original Hugging Face training split, while the test split is sampled from the original Hugging Face test split.
The resulting splits preserve the natural class imbalance as closely as possible: 10,000 training examples, 2,000 validation examples, and 3,000 test examples.
The toxic class represents 8\% of the train and validation splits and 7.90\% of the test split, as shown in Table~\ref{tab:natural-splits}.
This natural stratified split is the main evaluation setting.
The earlier 70/30 not-toxic/toxic split is kept as a controlled diagnostic setup, but it is no longer used as the main final evaluation split.
For the final distillation experiment, BERT-base teacher logits and probabilities were generated for the same natural train, validation, and test rows.
This keeps the supervised baselines, teacher model, and distilled BERT-tiny student aligned on the same data split.

\begin{table}[htbp]
\caption{Natural distribution stratified splits used for the main original-label experiments.}\label{tab:natural-splits}
\centering
\begin{tabular}{lrrrr}
\toprule
Split & Rows & Label 0 & Label 1 & Toxic \% \\
\midrule
Train & 10000 & 9200 & 800 & 8.0 \\
Validation & 2000 & 1840 & 160 & 8.0 \\
Test & 3000 & 2763 & 237 & 7.9 \\
\bottomrule
\end{tabular}
\end{table}

Preprocessing was intentionally minimal.
Handled empty values, leading and trailing spaces were removed, and repeated spaces were normalised. 
No lowercasing or punctuation removal was applied, because Transformer tokenizers are designed to process raw text.

\subsection{Evaluation metrics}

The main classification metrics are accuracy, precision, recall, F1 score, and the confusion matrix. 
Because the task is class imbalanced, F1 score is treated as more informative than accuracy measure.
In addition, model compactness and computational cost are evaluated using model size on disk, average inference time per example, and examples processed per second on CPU.

\subsection{Supervised baselines and BERT-base teacher}

The first baseline is a classical TF-IDF and Logistic Regression classifier.
It represents the vector space and linear classification part of the NLP course.
This baseline was retrained on the natural distribution split.
It was using unigram and bigram TF-IDF features, a maximum of 20,000 features, and balanced Logistic Regression.
The compact supervised student baseline is \texttt{prajjwal1/bert-tiny}.
It's finetuned on the same natural distribution original labels.
The local supervised teacher is a pretrained \texttt{bert-base-uncased} model.
It's finetuned on the same natural distribution original labels.
It is important to distinguish this from BERT pretraining.
The Transformer models are not trained from scratch, but finetuned for toxic comment classification.
Because the dataset is imbalanced, the BERT-tiny and BERT-base training runs use weighted cross-entropy with class weights computed from the training distribution.

For all supervised models, a default threshold of 0.5 was evaluated.
In addition, the range of thresholds from 0.05 to 0.95 was checked on the validation set.
And the threshold with the best validation F1 was selected.
The selected threshold was then evaluated once on the test set.
Threshold tuning is validation based and does not use test labels for model selection.

\subsection{Training dynamics and checkpoint selection}

Training histories were recorded for the BERT-base teacher and the soft-label distilled BERT-tiny student.
Figure~\ref{fig:bert-base-training-dynamics} shows that BERT-base reached its best validation F1 at epoch 3.
After this checkpoint, training loss continued decreasing while validation loss increased, suggesting overfitting.
Figure~\ref{fig:distilled-training-dynamics} shows that the distilled BERT-tiny model continued improving until epoch 5.
Test metrics are reported from the best validation checkpoint for each model.

\begin{figure}[htbp]
\centering
\begin{minipage}{0.48\textwidth}
\centering
\includegraphics[width=\linewidth]{../apprentice_model/results/figures/bert_base_learning_curve_f1.png}
\end{minipage}
\hfill
\begin{minipage}{0.48\textwidth}
\centering
\includegraphics[width=\linewidth]{../apprentice_model/results/figures/bert_base_learning_curve_loss.png}
\end{minipage}
\caption{BERT-base teacher validation F1 and loss curves.}
\label{fig:bert-base-training-dynamics}
\end{figure}

\begin{figure}[htbp]
\centering
\begin{minipage}{0.48\textwidth}
\centering
\includegraphics[width=\linewidth]{../apprentice_model/results/figures/bert_tiny_distilled_t2_a07_f1.png}
\end{minipage}
\hfill
\begin{minipage}{0.48\textwidth}
\centering
\includegraphics[width=\linewidth]{../apprentice_model/results/figures/bert_tiny_distilled_t2_a07_loss.png}
\end{minipage}
\caption{Soft-label distilled BERT-tiny validation scores and loss curves.}
\label{fig:distilled-training-dynamics}
\end{figure}

Table~\ref{tab:natural-baselines} shows the main comparable test results at the default threshold on the 3,000 example natural distribution test split.

\begin{table}[htbp]
\caption{Main supervised comparison on the natural distribution 3,000 example test split using default threshold 0.50.}\label{tab:natural-baselines}
\centering
\scriptsize
\setlength{\tabcolsep}{3pt}
\begin{tabular}{@{}lrrrrrrrr@{}}
\toprule
Model & Accuracy & Precision & Recall & F1 & TN & FP & FN & TP \\
\midrule
TF-IDF + Logistic Regression & 0.8967 & 0.3719 & 0.4473 & 0.4061 & 2584 & 179 & 131 & 106 \\
BERT-tiny supervised & 0.8863 & 0.3844 & 0.7300 & 0.5036 & 2486 & 277 & 64 & 173 \\
BERT-base supervised teacher & 0.9313 & 0.5508 & 0.7089 & 0.6199 & 2626 & 137 & 69 & 168 \\
\bottomrule
\end{tabular}
\end{table}

The BERT-base supervised teacher substantially improves over the TF-IDF baseline on the natural test split.
At the default threshold, BERT-base reaches an F1 score of 0.620, compared with 0.504 for BERT-tiny and 0.406 for TF-IDF.
BERT-tiny is lower than BERT-base, but it clearly improves recall and F1 over the classical baseline.
The TF-IDF model detects 106 of 237 toxic test examples at threshold 0.5.
BERT-tiny detects 173 of 237 toxic examples, but also produces 277 false positives.
This shows that BERT-tiny over-predicts the toxic class at the default threshold.
BERT-base detects 168 toxic examples while reducing false positives to 137, giving the best precision-recall balance.

The default-threshold test confusion matrix for the TF-IDF baseline is
\begin{equation}
\begin{bmatrix}
2584 & 179 \\
131 & 106
\end{bmatrix},
\end{equation}
the BERT-tiny supervised model obtains
\begin{equation}
\begin{bmatrix}
2486 & 277 \\
64 & 173
\end{bmatrix},
\end{equation}
and the BERT-base supervised teacher obtains
\begin{equation}
\begin{bmatrix}
2626 & 137 \\
69 & 168
\end{bmatrix}.
\end{equation}
These matrices show the tradeoff between the compact and larger Transformer models.
BERT-tiny reduces false negatives compared with TF-IDF, but increases false positives.
BERT-base gives the strongest overall result, reducing both false positives and false negatives compared with TF-IDF.
This justifies using BERT-base as the teacher model for later soft-label distillation into BERT-tiny.
The finetuned BERT-base model saved in \texttt{results/bert\_base\_supervised/} matches the best validation checkpoint, that was at epoch 3 with validation F1 equal to 0.615.
The model logits and probabilities were exported for the train, validation, and test splits and used for the soft-label distillation experiment.

The BERT-base validation curve is reported in Table~\ref{tab:bert-base-history}.
The model reaches its best validation F1 at epoch 3.
After that point, training loss continues to decrease but validation loss increases, which indicates mild overfitting.

\begin{table}[htbp]
\caption{BERT-base supervised teacher validation history.}\label{tab:bert-base-history}
\centering
\begin{tabular}{rrrrr}
\toprule
Epoch & Train loss & Validation loss & Validation accuracy & Validation F1 \\
\midrule
1 & 0.413 & 0.310 & 0.930 & 0.614 \\
2 & 0.212 & 0.322 & 0.863 & 0.500 \\
3 & 0.117 & 0.351 & 0.933 & 0.615 \\
4 & 0.069 & 0.638 & 0.928 & 0.598 \\
5 & 0.044 & 0.834 & 0.939 & 0.599 \\
\bottomrule
\end{tabular}
\end{table}

\begin{table}[htbp]
\caption{Validation threshold tuning summary for the natural distribution experiment.}\label{tab:threshold-tuning}
\centering
\small
\begin{tabular}{@{}lrrr@{}}
\toprule
Model & Tuned threshold & Test F1 at 0.50 & Test F1 tuned \\
\midrule
TF-IDF + Logistic Regression & 0.52 & 0.406 & 0.400 \\
BERT-tiny supervised & 0.75 & 0.504 & 0.543 \\
BERT-base supervised teacher & 0.49 & 0.620 & 0.618 \\
BERT-tiny distilled, $T=2$, $\alpha=0.7$ & 0.75 & 0.565 & 0.559 \\
\bottomrule
\end{tabular}
\end{table}

The validation-tuned results are treated as an additional calibration analysis.
Threshold tuning decreases test F1 for TF-IDF and BERT-base, but improves BERT-tiny from 0.504 to 0.543.
This is consistent with the default threshold behaviour of BERT-tiny.
Raising the threshold reduces false positives and improves its test F1.
For the main comparison, however, Table~\ref{tab:natural-baselines} keeps the default threshold fixed at 0.50 for all models.
For the soft label distilled BERT-tiny model, validation tuning selected threshold 0.75.
It is corresponding to the tuned test F1 of 0.5590.
The default threshold test F1 was 0.5649.
Therefore, the default threshold 0.50 remains the main comparable result for the distilled model as well.

\subsection{BERT-base teacher analysis}

The final teacher model is the finetuned \texttt{bert base uncased} classifier trained on the natural distribution training split.
It is a stronger supervised model whose logits and probabilities are used as the distillation signal for BERT-tiny.
Teacher logits were exported for the same 10,000 training, 2,000 validation, and 3,000 test examples used in the main experiment.
This makes the teacher and student predictions directly comparable.

\begin{table}[htbp]
\caption{BERT-base teacher performance on the natural test split at the default threshold 0.50.}
\label{tab:bert-base-teacher-summary}
\centering
\begin{tabular}{lrrrrrrrr}
\toprule
Model & Accuracy & Precision & Recall & F1 & TN & FP & FN & TP \\
\midrule
BERT-base teacher & 0.931 & 0.551 & 0.709 & 0.619 & 2626 & 137 & 69 & 168 \\
\bottomrule
\end{tabular}
\end{table}

The teacher has the strongest default threshold F1 among the supervised models.
It also provides the probability distribution over the two classes, not only a hard class prediction.
The distilled BERT-tiny student is trained to preserve information from the teacher's softened output distribution.
It uses the original dataset labels in the cross entropy part of the loss.
The teacher should not be treated as ground truth.
It is a stronger model used as a learning signal, and its errors or biases can still be transferred to the student.

\subsection{Distilled student results}

The first tries of distillation used hard teacher labels.
This was a useful starting point because it made the pipeline simple.
The student was trained as a normal classifier, but with teacher predictions as targets.
However, hard labels don't use the information about teacher uncertainty.
For example, a teacher prediction with probabilities 0.51 to 0.49 is treated the same as a prediction with probabilities 0.99 to 0.01.
For this reason, the final experiment switched to softlabel temperature distillation using the BERT-base teacher logits.
This allows the BERT-tiny student to learn from the teacher's full output distribution rather than only from the teacher's argmax prediction.

\begin{table}[htbp]
\caption{Soft label distilled BERT-tiny result on the natural test split.}\label{tab:distilled-results}
\centering
\begin{tabular}{lrrrr}
\toprule
Model & Accuracy & Precision & Recall & F1 \\
\midrule
BERT-tiny distilled, $T=2$, $\alpha=0.7$ & 0.9117 & 0.4624 & 0.7257 & 0.5649 \\
\bottomrule
\end{tabular}
\end{table}

The main final experiment uses softlabel temperature distillation from the local BERT-base teacher into BERT-tiny on the natural stratified split.
In this experiment, the teacher is the finetuned BERT-base model.
The student is \texttt{prajjwal1/bert-tiny}, the temperature is $T=2.0$, and $\alpha=0.7$.
The loss is
\begin{equation}
    \mathcal{L}
    =
    (1-\alpha)\,\mathcal{L}_{CE}
    +
    \alpha T^2\,\mathcal{L}_{KL}.
\end{equation}
At the default threshold of 0.50, the softlabel distilled BERT-tiny model reaches test F1 0.565 on the natural split.
This improves over the supervised BERT-tiny baseline by 0.061 F1.
It also recovers 52.65\% of the F1 difference between supervised BERT-tiny and BERT-base.
Compared with the BERT-base teacher's hard predictions, the distilled student got test F1 0.679 and agreement 0.928.
The main effect is that distillation reduces false positives compared with supervised BERT-tiny.
And it preserves similar recall.

\subsection{Error analysis candidates}

To perform manual error analysis, I compared the original label, BERT-base teacher prediction, supervised BERT-tiny prediction, and distilled BERT tiny prediction on the natural test split.
The aggregate counts show that softlabel distillation mostly improves BERT-tiny by reducing false positives.
It fixes 99 supervised BERT-tiny false positives and adds 22 new false positives.
So the net reduction is 77 false positives.
False negatives stay almost unchanged.
11 are fixed and 12 are worsened.
This matches the confusion matrices.
Supervised BERT-tiny has FP=277 and FN=64, and distilled BERT-tiny has FP=200 and FN=65.
Table~\ref{tab:error-analysis-examples} shows six chosen examples for manual inspection.

\begin{table}[htbp]
\caption{Curated error-analysis examples from the natural test split.}
\label{tab:error-analysis-examples}
\centering
\tiny
\setlength{\tabcolsep}{1.5pt}
\begin{tabular}{p{1.8cm}p{4.0cm}rrrrp{3.0cm}}
\toprule
Case & Text excerpt & Gold & Teacher & Sup. & Dist. & Short interpretation \\
\midrule
Fixed FP & It is good that the world starts to understand the racist mentality of Isr*el... & 0 & 0 & 1 & 0 & Distillation lowers toxicity probability. \\
Fixed FP & For a laugh, try Infowars. It makes Breitbart look the venerable Times of London. & 0 & 0 & 1 & 0 & Media criticism stays non-toxic. \\
Fixed FN & People who den*ate any American's right to protest are nothing more than cretinous boors. & 1 & 1 & 0 & 1 & Insulting wording crosses threshold. \\
Worse FP & I would hate to be the one who had to read those letters. Prayers for all those children... & 0 & 0 & 0 & 1 & Negative text crosses threshold. \\
Still FP & We do not have a President. The guy does not need any help tarnishing himself... & 0 & 0 & 1 & 1 & Harsh political criticism is flagged. \\
Still FN & Two-timer. Cad. & 1 & 0 & 0 & 0 & Very short toxic-labeled phrase is missed. \\
\bottomrule
\end{tabular}
\end{table}

\subsection{Token saliency examples}

For the quality validation i computed the saliency scores.
For each token, I measure how strongly the toxic class score depends on that token’s embedding.
Tokens with higher scores are treated as more influential for the toxic prediction.
The scores are normalized per example and model, so they are illustrative only.
Tokens with high score should be interpreted as tokens with high saliency for the toxic class score.
Table~\ref{tab:saliency-examples} shows four compact examples.

\begin{table}[htbp]
\caption{Qualitative token saliency examples for supervised and distilled BERT-tiny.}
\label{tab:saliency-examples}
\centering
\tiny
\setlength{\tabcolsep}{1pt}
\begin{tabular}{p{0.45cm}p{1.25cm}p{0.35cm}p{0.55cm}p{0.55cm}p{2.2cm}p{2.2cm}p{2.2cm}}
\toprule
Row & Case & Y & Sup. & Dist. & Sup. tokens & Dist. tokens & Short note \\
\midrule
67 & Fixed FP & 0 & 0.791 & 0.259 & racist, mentality, starts & racist, mentality, palestinian & Loaded tokens remain salient; distilled probability is lower. \\
79 & Fixed FP & 0 & 0.547 & 0.020 & laugh, Infowars, Breitbart & laugh, Infowars, Breitbart & Media tokens stay salient; distilled probability is lower. \\
69 & Fixed FN & 1 & 0.331 & 0.527 & boors, unpatriotic, cretinous & unpatriotic, boors, cretinous & Distillation raises probability for insulting wording. \\
46 & Still FP & 0 & 0.623 & 0.649 & hell, bizarreness, guy & hell, bizarreness, indefensible & Harsh political criticism remains salient. \\
\bottomrule
\end{tabular}
\end{table}

\subsection{Minimal pair bias check}

A small minimal pair set was manually constructed.
It was used to compare neutral identity mentions, non identity negative statements, direct insults, and identity attacks.
This is a qualitative check, not a full fairness audit.
The aggregate results are shown in Table~\ref{tab:minimal-pair-bias}.
Counts show correct predictions over manually assigned expected labels.

\begin{table}[htbp]
\caption{Minimal pair qualitative check by category}
\label{tab:minimal-pair-bias}
\centering
\tiny
\begin{tabular}{lcccc}
\toprule
Model & Neutral ID & Negative & Insult & ID attack \\
\midrule
BERT-base & 5/7 & 3/3 & 3/3 & 5/5 \\
BERT-tiny sup. & 6/7 & 3/3 & 3/3 & 5/5 \\
BERT-tiny dist. & 6/7 & 3/3 & 3/3 & 5/5 \\
\bottomrule
\end{tabular}
\end{table}

All models correctly classified the simple direct insults and identity attacks in this small set.
However, neutral identity mentions were not always classified correctly.
In particular, ``I am gay.'' was classified as toxic by all three models, and ``I am Muslim.'' was classified as toxic by the BERT-base teacher.
This suggests possible sensitivity to identity terms, and distillation does not guarantee removal of teacher or dataset biases.

\subsection{Efficiency benchmark}

Efficiency was measured on CPU using the test split.
The benchmark includes model size on disk, average inference time per example, and throughput in examples per second.
For the natural split comparison, I benchmarked all final models on CPU over the same 3,000 test examples with one repeat.
For Transformer models, the main comparison uses batch size 16 and maximum sequence length 256.
The results are shown in Table~\ref{tab:september-cpu-benchmark}.

\begin{table}[htbp]
\caption{CPU benchmark on the natural distribution test split using batch size 16 for Transformer models.}
\label{tab:september-cpu-benchmark}
\centering
\scriptsize
\begin{tabular}{lrrrr}
\toprule
Model & F1 & Size (MB) & ms/example & Examples/sec \\
\midrule
TF-IDF + LR & 0.4061 & 0.89 & 0.1333 & 7499.06 \\
BERT-tiny supervised & 0.5036 & 17.42 & 5.2392 & 190.87 \\
BERT-tiny distilled & 0.5649 & 17.42 & 5.4296 & 184.18 \\
BERT-base teacher & 0.6199 & 418.35 & 288.4259 & 3.47 \\
\bottomrule
\end{tabular}
\end{table}

BERT-base gives the best F1, but it is much larger and slower than the compact models.
The soft label distilled BERT-tiny improves over supervised BERT-tiny while keeping the same compact model size and similar inference speed.
Compared with BERT-base, distilled BERT-tiny is about 24 times smaller and about 53 times faster on CPU at batch size 16.
This shows how the distillation can be useful.
It is a better compact model with much lower inference cost than the teacher.
The exact speed numbers are hardware-dependent.

\section{Concluding Remarks}\label{sec:conclusion}

In this project a complete apprentice model pipeline was implemented and evaluated for toxic comment classification.
The results show that a compact BERT-tiny student can be adapted to imitate a stronger BERT-base teacher on a narrow NLP task.
On the natural distribution original-label experiment, the classical TF-IDF and Logistic Regression baseline reached an F1 score of 0.406 on the test split at the default threshold.
The supervised BERT-tiny baseline reached an F1 score of 0.504.
The finetuned BERT-base supervised teacher reached an F1 score of 0.62 on the same test split.
This confirms that contextual Transformer representations provide a high improvement over the classical sparse vector baseline for this imbalanced toxicity task.
The final soft label distilled BERT-tiny model reached F1 0.565 against original labels at the default threshold.
This improves over supervised BERT-tiny by 0.061 F1 and recovers 52.65\% of the gap between supervised BERT-tiny and the BERT-base teacher.
Against BERT-base teacher hard predictions, the distilled model reached F1 0.679 and agreement 0.928, showing partial transfer of the teacher's decision behaviour.

The experiments also highlight important tradeoffs.
The TF-IDF baseline is extremely efficient and interpretable, but weaker in F1 score.
BERT-tiny and BERT-base provide better classification quality, but are significantly larger and slower.
BERT-base gives the best F1, but the CPU benchmark shows that it is much larger and slower than BERT-tiny.
At batch size 16, distilled BERT-tiny is about 24 times smaller and about 53 times faster than BERT-base.
It is improving over the supervised BERT-tiny baseline.
Distillation does not change model size or speed because the architecture remains the same, but it changes the decision boundary of the student.
This supports the idea that knowledge distillation can be used not only for compression, but also for task specific behavioural adaptation.

There are several limitations.
First, the BERT-base teacher is not an objective ground truth and may transfer errors or biases to the student.
Second, the minimal pair check suggests that all models can still be sensitive to identity terms.
Third, only one soft label distillation setting, $T=2.0$ and $\alpha=0.7$, is used as the accepted final run.
Finally, validation threshold tuning does not always improve test F1.
This means that threshold calibration should be interpreted carefully.

Future work could extend the project in several directions.
The most direct extension is to compare more temperatures, alpha values, and calibration methods for the BERT-base to BERT-tiny distillation setup.
Additional work could compare multiple teachers, such as an OpenAI LLM and an open-source toxicity classifier.
Or include explanation based distillation.
Finally, a larger and more systematic bias evaluation would make the qualitative minimal-pair findings more reliable.
The CPU efficiency benchmark and model size comparison for the BERT-base to BERT-tiny soft label distillation experiment should mainly be interpreted as a relative comparison between models under the same setup.

\begin{thebibliography}{9}

\bibitem{hinton2015distilling}
Hinton, G., Vinyals, O., Dean, J.: Distilling the knowledge in a neural network. arXiv preprint arXiv:1503.02531 (2015)

\bibitem{devlin2019bert}
Devlin, J., Chang, M.-W., Lee, K., Toutanova, K.: BERT: Pre-training of deep bidirectional transformers for language understanding. In: Proceedings of NAACL-HLT, pp. 4171--4186 (2019)

\bibitem{jiao2020tinybert}
Jiao, X., Yin, Y., Shang, L., Jiang, X., Chen, X., Li, L., Wang, F., Liu, Q.: TinyBERT: Distilling BERT for natural language understanding. In: Findings of the Association for Computational Linguistics: EMNLP 2020, pp. 4163--4174 (2020)

\end{thebibliography}

\end{document}
